\documentclass[../../../include/open-logic-section]{subfiles}

\begin{document}

\olfileid[pt]{sth}{story}{predicative}	

\olsection{Predicativo e Impredicativo}

O conjunto de Russell, $R$, foi definido via $\Setabs{x}{x \notin
x}$. Dito com mais detalhe, $R$ seria o conjunto que contém todos
e apenas aqueles conjuntos que não são membros de si próprios. Assim, ao definir
$R$, quantificamos sobre o domínio que conteria $R$ (se este
existisse).

Isto é uma definição \emph{impredicativa}. Mais em geral, podemos
dizer que uma definição é impredicativa sse quantifica sobre um domínio
que contém o objeto que está a ser definido.  
	
Na sequência dos paradoxos, Whitehead, Russell, Poincar\'{e} e
Weyl rejeitaram tais definições impredicativas como ``viciosamente
circulares'':
\begin{quote}
	Uma análise dos paradoxos a evitar mostra que todos eles
	resultam de uma espécie de círculo vicioso. Os círculos viciosos em
	questão surgem ao supor que uma coleção de objetos pode
	conter membros que só podem ser definidos por meio da
	coleção como um todo[\ldots. \textparagraph]

	O princípio que nos permite evitar totalidades ilegítimas
	pode enunciar-se do seguinte modo: ``Tudo o que envolve \emph{todos} os
	membros de uma coleção não pode ser um dos membros da coleção''; ou, inversamente:
	``Se, supondo que uma certa coleção tivesse um todo, ela teria
	membros só definíveis em termos desse todo, então a dita
	coleção não tem todo.'' Chamaremos a isto o ``princípio do círculo vicioso'',
	porque nos permite evitar os círculos viciosos
	envolvidos na suposição de totalidades ilegítimas.
	\citep[p.~37]{WhiteheadRussell1910}
\end{quote}
Se os acompanharmos na rejeição das definições impredicativas motivada pelo \emph{princípio do círculo vicioso},
então podemos tentar substituir o desastroso esquema de compreensão ingênua (de \olref[sth][story][rus]{sec}) por algo deste gênero: 

\begin{defish}
\emph{Compreensão predicativa.} Para toda a fórmula $\phi$ que quantifica apenas sobre conjuntos: existe o conjunto$^\prime$ $\Setabs{x}{\phi(x)}$.
\end{defish}

Enquanto os conjuntos$^{\prime}$ não forem conjuntos, não se seguirá contradição.  

Infelizmente, a compreensão predicativa não é muito
\emph{abrangente}. Afinal, introduz-nos novas entidades,
conjuntos$^\prime$. Logo teremos de considerar fórmulas que quantificam
sobre conjuntos$^\prime$. Se estas derem sempre um conjunto$^\prime$, então
o paradoxo de Russell surgirá de novo, bastando considerar o
conjunto$^\prime$ de todos os conjuntos$^\prime$ que não são membros de si próprios. Assim, a seguir o
mesmo fio de ideias, temos de dizer que uma fórmula que quantifica sobre
conjuntos$^\prime$ produz um conjunto$^{\prime\prime}$ correspondente. E então
precisaremos de conjuntos$^{\prime\prime\prime}$,
conjuntos$^{\prime\prime\prime\prime}$, etc. Para evitar uma enxurrada de primos,
será mais fácil pensar nestes como conjuntos$_0$, conjuntos$_1$, conjuntos$_2$,
conjuntos$_3$, conjuntos$_4$,\ldots. E isto dar-nos-ia entrada na
teoria (simples) dos tipos. 

Há algumas objeções óbvias a tal teoria (embora não seja
óbvio que sejam objeções \emph{esmagadoras}). Em resumo:
a teoria resultante é pouco manejável; é prolixa ao
postular diferentes tipos de objetos; e não está claro, no
fim, que as definições impredicativas sejam sequer \emph{assim tão más}. 
	
Para ilustrar o último ponto, consideremos esta observação de
\citeauthor{Ramsey1925}:
\begin{quote}
	podemos referir-nos a um homem como o mais alto de um grupo, identificando-o
	por meio de uma totalidade da qual ele próprio é membro sem
	que haja círculo vicioso. \citep{Ramsey1925}
\end{quote}
O ponto de Ramsey é que ``o homem mais alto do grupo'' \emph{é} uma
definição impredicativa; mas é obviamente perfeitamente legítima. 

Alguém poderá responder que, neste caso, poderíamos identificar a pessoa mais alta por meios \emph{predicativos}. Por exemplo, talvez pudéssemos simplesmente
apontar para o homem em questão. A objeção às definições impredicativas,
então, teria claramente de se limitar a entidades que
\emph{só} possam ser identificadas impredicativamente. Mas mesmo assim, seria preciso ouvir mais, sobre por que tal ``impredicatividade essencial'' seria
tão grave.\footnote{Para mais, ver \citet{Linnebo2010}.}

Admitidamente, as definições impredicativas são péssimas notícias, se
quisermos que as definições nos forneçam algo como uma receita para
\emph{criar} um objeto. Pois, dada uma definição impredicativa, estar-se-ia
genuinamente preso a um círculo vicioso: para criar o
objeto especificado impredicativamente, seria preciso \emph{primeiro}
criar todos os objetos (incluindo o objeto especificado impredicativamente), uma vez que este é especificado em
termos de todos os objetos; logo seria preciso criar o
objeto especificado impredicativamente antes de o ter criado. Mas,
de novo, isto só é uma objeção séria a conjuntos especificados de modo ``essencialmente
impredicativo'', se pensarmos nos conjuntos como coisas que
\emph{criamos}. E nós (provavelmente) não os criamos.

Assim---para o bem ou para o mal---a abordagem que se tornou comum não
envolve tomar uma posição rígida sobre (im)\-pre\-di\-ca\-ti\-vi\-da\-de.
Envolve antes o que hoje se considera a abordagem
cumulativa-iterativa. No fim, isto permitir-nos-á estratificar os conjuntos em
``estágios''---um \emph{pouco} como a abordagem predicativa estratifica
entidades em conjuntos$_0$, conjuntos$_1$, conjuntos$_2$, \ldots---mas não
postularemos diferença de natureza entre eles. 

\end{document}
